\begin{thm}[not all lines wrap/acyclic cases]\label{not-all-wrap}
The following are equivalent:
\begin{enumerate}\itemsep=0pt
\item[$1.$] $U^{\trop}$ contains a {line} which does not wrap.
\item[$2.$] Some compactification of $U$ admits a toric model $Y\rar \?{Y}$ for which all the non-toric blowups are on divisors corresponding to rays in one half of $N_{\?{Y}}$. I.e., there is some seed~$S$ for which all of the non-frozen vectors' images in~$p_2^*(N)$ lie in one half of the plane.
\item[$3.$] Cluster varieties corresponding to~$U$ are {acyclic}.
\item[$4.$] The intersection of the Langlands dual {cluster complex} $\s{C}\subset \s{X}^{\trop}$ with $U^{\trop}$ is non\-empty.
\item[$5.$] There exists a {global monomial} on~$U$.
\item[$6.$] The {quadratic form} $Q$ on $D^{\perp}$ is negative definite, and $Q|_{D^{\perp}_{\Eff}}$ is a direct sum of $A_{n_i}$'s. In fact, it is $A_{d'_1-1}\oplus \cdots \oplus A_{d'_m-1}$, where the $(d'_i)$'s are the modified multipliers for a coprime seed corresponding to~$U$ $($equivalently, $d_i'$ is the number of non-toric blowups on~$D_i$ in a~toric model for a~compactification~$U)$.
\end{enumerate}
\end{thm}
\begin{proof}(1)$\Leftrightarrow$(2) is Lemma \ref{model}. (2)$\Leftrightarrow$(3) was observed in Section~\ref{quivers}.

For (2)$\Leftrightarrow$(4), note that for some seed vector $e_i$ for a seed $S$, the set $\{e_i\geq 0\}\cap U^{\trop}$ is the same as the set $(v_i\wedge \cdot) \geq 0$, where $\wedge$ is the symplectic form on $U^{\trop}$ induced by $[ \cdot,\cdot]$. The intersection of these positive half-spaces for all non-frozen $e_i$'s is clearly nonempty if and only if~$S$ is as in~(2).

 (1)$\Rightarrow$(5) follows from Lemma \ref{GlobalMonomial}. For (5)$\Rightarrow$(1), note that for a global monomial $\vartheta_q$, the tropicalization $\vartheta_q^{\trop}$ is positive somewhere, and so Lemma \ref{NegativeFibers} implies that the fibers $\vartheta_q^{\trop}=d<0$ are lines which do not wrap.

 (6)$\Rightarrow$(1) because if every line does wrap (possibly infinitely many times), then we have seen that either $Q$ is not negative-definite or $Q|_{D^{\perp}_{\Eff}}$ is of type $D_n$ or $E_n$.

 For (2)$\Rightarrow$(6), first note that $Q$ is negative definite on $D^{\perp}$ by positivity of $U$. Now, let $(Y,D)\rar \big(\?{Y},\?{D}\big)$ be the toric model corresponding to a seed with all non-toric blowups corresponding to rays in one half of the plane $N_{\?{Y}}$. For any curve $\?{C}$ in $\?{Y}$, $\sum \big(\?{C}\cdot \?{D}_i\big) v_i=0$ where $v_i$ is the primitive vector in $N_{\?{Y}}$ corresponding to $\?{D}_i$. If $\?{C}$ is the image of an irreducible effective curve $C\in D^{\perp}$, then $\?{C}\cdot \?{D}_i \geq 0$ for all $i$, and $\?{C}\cdot \?{D}_i$ can only be positive if there is a non-toric blowup point somewhere in $\?{C} \cap \?{D}_i$. Thus, each $\?{C}\cdot \?{D}_i$ must actually be $0$, so $C$ must have been supported on an exceptional divisor. Thus, $D^{\perp}_{\Eff}$ is generated by classes obtained by taking the $d_i'$ blowups to be infinitely near, and then taking the $d_i'-1$ exceptional divisors which do not intersect $D$.
\end{proof}

Let $\s{C}_U$ denote the union of all cones $\sigma_L$ for lines $L$ which do not wrap, where $\sigma_L$ is defined as in Lemma~\ref{model}. We note that the argument for (2)$\Leftrightarrow$(4) above can be modified to prove the following:
\begin{prop} $\s{C}_U$ is the intersection of the Langlands dual cluster complex $\s{C}$ with $U^{\trop}\subset \s{X}^{\trop}$.
\end{prop}
This justifies \cite{Man1} calling $\s{C}_U$ the cluster complex.

\begin{thm}[no lines wrap/finite-type cases]\label{no-wrap}
The following are equivalent:
\begin{enumerate}\itemsep=0pt
\item[$1.$] No {lines} in $U^{\trop}$ wrap.
\item[$2.$] No sheet of the {developing map} is convex.
\item[$3.$] The Laurent phenomenon holds for the $\s{X}$-space, meaning that each $X_i$ is a global monomial. Furthermore, the {global monomials} form an additive basis for the global function on~$U$.
\item[$4.$] The inverse {monodromy} matrix $\mu^{-1}$ is a Kodaira matrix of type $I_k$, $II$, $III$, or~$IV$.
\item[$5.$] Cluster structures for $U$ are of {finite type}, meaning that they have only a finite number of distinct seeds.
\item[$6.$] For some equivalent maximally factored seed, the corresponding quiver $($after removing frozen vectors$)$ is of type $A_1^k$ $(k\in \bb{Z}_{\geq 0})$, $A_2$, $A_3$, or $D_4$.
\item[$7.$] The Langlands dual {cluster complex} $\s{C}\subseteq \s{X}^{\trop}$ contains all of~$U^{\trop}$, and in fact is all of~$\s{X}^{\trop}$.
\end{enumerate}
\end{thm}
\begin{proof}(1)$\Leftrightarrow$(2) is obvious. (1)$\Leftrightarrow$(3) follows from Lemma~\ref{GlobalMonomial}.

To see that (1) implies (5), we need Lemma~\ref{model}, which says that for any line~$L_q^{d<0}$ which does not wrap, there are only finitely many $(-1)$-curves hitting boundary divisors corresponding to rays in the cone $\sigma_{L}$ bounded by $L_q^{d<0}(\pm \infty)$. Since no lines wrap, we can cover $U^{\trop}$ by finitely many cones of the form $\sigma_{L}$, and so there are only finitely many $(-1)$-curves in $Y$ hitting the boundary. Since seeds correspond to certain finite subsets of this collection of $(-1)$-curves, the claim follows.

(5)$\Leftrightarrow$(6) follows from a well-known result of \cite{FZ2}, which says that a cluster algebra is of finite type if and only if the matrix $(-|\epsilon_{ij}|+2\delta_{ij})_{i,j\in I\setminus F}$ is a finite type Cartan matrix. One easily checks that the only quivers of this type which produce rank $2$ cluster varieties are those listed in the statement of theorem, along with types $B_2$, $B_3$, $C_3$, and $G_2$, which are equivalent to types~$A_3$,~$D_4$,~$D_4$, and $D_4$ again, respectively, in the sense of Definition \ref{EquivalentClusters}.

 One can easily check (6)$\Rightarrow$(4) by explicit computations: the $A_1^k$, $A_2$, $A_3$, and $D_4$ quivers correspond to the $I_k$, $II$, $III$, and $IV$ matrices, respectively. (4)$\Rightarrow$(1) is now automatic.

For (5)$\Leftrightarrow$(7), recall that seeds are in bijection with cones of the cluster complex. For any boundary wall $W$ of any cone in $\s{C}$, both sides of $W$ will always be in $\s{C}$, so if there are only finitely many cones, then~$\s{C}$ must fill up all of~$\s{X}^{\trop}$. Conversely, if there are infinitely many cones, then they must ``bunch up'' near some ray $\rho$ which is not in $\s{C}$.
\end{proof}

Table \ref{tab:no wrap} lists the cases where no lines wrap, along with their basic properties. We once again use the notation $(d_1,d_2,d_3)$ to indicate that such a Looijenga pair can be obtained by starting with the toric variety $\big(\bb{P}^2,D=D_1+D_2+D_3\big)$, and then blowing up $d_1$, $d_2$, and $d_3$ points on~$D_1$,~$D_2$, and~$D_3$, respectively.
\begin{table}[h]\centering
 \begin{tabular}{ | l | l | l | l | l|}
 \hline
 Quiver & Kodaira matrix & Cartan form $Q$ & Monodromy $\mu$ & $(d_1,d_2,d_3)$ \\ \hline
 $A_1^k$ ($k\geq 0$) & $I_{k}$ & $A_{k-1}$ & 	 $\left(\begin{matrix}
		1 &-k \\
	 0 &1
		\end{matrix}\right)$				 		& $(k,0,0)$ \\ \hline
 $A_2$ & $II$ & $A_0$ & 		$\left(\begin{matrix}
		0 &-1 \\
	 1 &1
		\end{matrix}\right)$		& $(1,1,0)$	 \\ \hline
 $A_3$ & $III$ &$A_1$ &	$\left(\begin{matrix}
		0 &-1 \\
	 1 &0
		\end{matrix}\right)$	 	 & $(2,1,0)$	 \\ \hline
 $D_4$ & $IV$ &$A_2$ & 	$\left(\begin{matrix}
		-1 &-1 \\
	 1 &0
		\end{matrix}\right)$			& $(3,1,0)$			 \\ \hline
 \end{tabular}
 \caption{Cases where no lines wrap.}\label{tab:no wrap}
\end{table}

\begin{rmk}\label{NeedFrozen}
Without frozen vectors, the $I_k$ cases, $k\geq 0$, are actually of rank $0$. Thus, although we tend to ignore frozen vectors, they are necessary for constructing these examples. They are also necessary for many other examples~-- this was reflected in Construction \ref{LooijengaConstr} when we required that the vectors $u_1,\dots,u_m$ generate $N_{\?{Y}}$.
\end{rmk}

\begin{prop}[some lines wrap and some do not]\label{SomeWrap} The following are equivalent:
\begin{enumerate}\itemsep=0pt
\item[$1.$] Some {lines} in $U^{\trop}$ wrap, while others do not.
\item[$2.$] Some $($but not all$)$ sheets of the {developing map} are convex.
\item[$3.$] Cluster varieties corresponding to~$U$ are {acyclic} but not of {finite type}.
\item[$4.$] The {monodromy} satisfies $\Tr(\mu) \leq -2$, and if there is equality, then $\mu$ is conjugate to $
	\left(\begin{smallmatrix}
		-1 &a \\
		0 &-1
		\end{smallmatrix}\right)$ for some $a<0$.
\end{enumerate}
\end{prop}
\begin{proof}(1)$\Leftrightarrow$(2) is easy, and (1)$\Leftrightarrow$(3) follows immediately from Theorems~\ref{not-all-wrap} and~\ref{no-wrap}. We see the equivalence of (1) with (4) because Theorems \ref{NegD}, \ref{NegSD}, \ref{all-wrap}, and \ref{no-wrap} have shown that all other possibilities for $\mu$ are equivalent to other cases.
\end{proof}

